\section{Scope and the exact logical claim}
The objective is an explicitly specified fixed polynomial with integer
coefficients whose zero existence, in the domain stated for the construction,
is equivalent to existence of a finite syntactic refutation of ZFC.
Writing $D$ for that domain, the ordinary arithmetic equivalence has the form
\[
 \mathsf{ZFC}\vdash
 \bigl(\neg\exists\mathbf x\in D^n\;P(\mathbf x)=0\bigr)
 \longleftrightarrow\operatorname{Con}(\mathsf{ZFC}).
\]
The equivalence does not prove either side. The fixed recursive presentation
and proof coding are those of the inherited source calibration. No consistency
assumption is used to choose coefficients or establish the zero-set reductions.

The new proofs are ordinary unbounded mathematics relative to the retained
source theorems. Neither a successful program replay nor source checksums
supply a fresh end-to-end Lean, Rocq, or object-language ZFC certificate.
The inherited compiler, Pell and formal-system calibration interfaces remain
substantive mathematical dependencies. Exact finite computation checks the
new transcription, identities, degree certificates and data integrity; it does
not replace those unbounded hypotheses.

The domains $\N$, $\Z_{>0}$ and $\Z$ are distinct comparison classes.
Each construction states its domain explicitly. A named polynomial register
is an abbreviation to be substituted, not an additional existential unknown.
Proof-only divisions, Pell sequences and radicals are not operations of the
supplied division-free equations.

Circuit cost counts live binary additions, subtractions and multiplications;
a square or multiplication by a fixed numeral is one multiplication.
An exact degree is distinguished from a syntactic DAG upper bound. Fixed
coefficient bit length, variable count, total degree, circuit cost, support
and coefficient height are separate criteria. A gain in one criterion need
not dominate another construction.
